% =============================================================================
% Assignments/CS301/10-file-organization-hashing-coding-assignment.tex
% Coding Assignment 10: Hash Tables -- Chaining and Linear Probing
% (Week 10, syllabi/CS301.md Week 10 lab: "Hash-table construction and
% collision handling", instructor-supplied)
%
% Student-facing problem statement. Instructor-only files live in
% Assignments/CS301/10-file-organization-hashing-coding/:
%   reference_solution.c  (NEVER published, gitignored)
%   test_hidden.c         (NEVER published, gitignored)
%   autograder.sh         (NEVER published, gitignored)
%   hash.h, starter_hash.c, test_visible.c  (provided to students)
%
% Fresh instances: the deck's own traces (m=11 with 12,44,13,88,23,94;
% h("cat")=1 mod 97; the 2,000-identifier table) are NOT reused. The worked
% examples here use m=5 chaining over alpha/beta/gamma/delta and m=11 linear
% probing over the same four strings, with h("dog")=58 mod 101.
%
% Compile with (macOS/Linux): cd Assignments/CS301 &&
%   TEXINPUTS=../../Preambles:$TEXINPUTS xelatex -interaction=nonstopmode 10-file-organization-hashing-coding-assignment.tex (3 passes)
% =============================================================================

\documentclass[11pt]{article}
\usepackage[margin=1in]{geometry}
% =============================================================================
% Preambles/header.tex — shared LaTeX/Beamer preamble for this project
%
% Use from a Beamer lecture by prepending:
%     \input{header}
% and compiling with TEXINPUTS=../Preambles:$TEXINPUTS (the /compile-latex
% skill does this automatically).
%
% PALETTE CONTRACT. Color names defined here MUST match the names in
% Quarto/theme-template.scss so Beamer and Quarto renderings use the same
% palette. The scripts/check-palette-sync.sh script enforces this.
%
% To customize: edit the HEX values below AND the matching $variable in
% Quarto/theme-template.scss, then run scripts/check-palette-sync.sh.
% =============================================================================

% -----------------------------------------------------------------------------
% Engine detection + encoding
%
% Our /compile-latex skill uses XeLaTeX, where inputenc is unnecessary and
% can produce encoding warnings. Only load inputenc under pdfTeX.
% -----------------------------------------------------------------------------
\usepackage{iftex}
\ifPDFTeX
  \usepackage[utf8]{inputenc}
\fi

\usepackage{amsmath, amssymb, amsthm}
\usepackage{booktabs}
\usepackage{graphicx}

% -----------------------------------------------------------------------------
% PALETTE — mirrors Quarto/theme-template.scss variable names.
% Load xcolor and define colors BEFORE anything that references them
% (notably hyperref, hypersetup, and the Beamer color assignments).
% -----------------------------------------------------------------------------
\usepackage{xcolor}

\definecolor{primary-blue}{HTML}{012169}      % headings, accents
\definecolor{primary-gold}{HTML}{B9975B}      % emphasis, borders
\definecolor{highlight-yellow}{HTML}{F2A900}  % markers, alerts
\definecolor{light-bg}{HTML}{E8EDF5}          % title / section backgrounds
\definecolor{jet}{HTML}{1A1A1A}               % body text

% Semantic colors (used in diagrams and annotations)
\definecolor{positive}{HTML}{15803D}          % good / observed / identified
\definecolor{negative}{HTML}{B91C1C}          % bad / problematic / bias
\definecolor{neutral}{HTML}{525252}           % reference / context

% Highlight accents (match .hi-* classes in the SCSS)
\definecolor{hi-slate}{HTML}{314F4F}
\definecolor{hi-green}{HTML}{2E8B57}
\definecolor{hi-red}{HTML}{C41E3A}

% Semantic aliases so skill templates and snippets can write
% \textcolor{accent}{...} without hard-coding "primary-blue".
\colorlet{accent}{primary-blue}
\colorlet{accent2}{primary-gold}

% -----------------------------------------------------------------------------
% Hyperref — loaded AFTER xcolor + palette so link colors resolve.
% -----------------------------------------------------------------------------
\usepackage{hyperref}
\hypersetup{colorlinks=true, linkcolor=primary-blue, urlcolor=primary-blue,
            citecolor=primary-blue, pdfborder={0 0 0}}

% -----------------------------------------------------------------------------
% Course code — set once per deck (after \input{header}, before \begin{document})
% via \coursecode{CS401}. Shown in the footer on every slide so a deck's course
% is identifiable without opening the title page. Empty by default.
% -----------------------------------------------------------------------------
\makeatletter
\newcommand{\coursecode}[1]{\gdef\@coursecodetext{#1}}
\gdef\@coursecodetext{}
\makeatother

% -----------------------------------------------------------------------------
% Beamer theme assignments (only apply under Beamer).
%
% We detect Beamer via \@ifclassloaded{beamer}, which is the documented,
% non-internal way. This must be wrapped in \makeatletter/\makeatother
% because \@ifclassloaded contains an @-catcode character.
% -----------------------------------------------------------------------------
\makeatletter
\@ifclassloaded{beamer}{%
  \usefonttheme{professionalfonts}
  \setbeamercolor{structure}{fg=primary-blue}
  \setbeamercolor{frametitle}{fg=primary-blue}
  \setbeamercolor{title}{fg=primary-blue}
  \setbeamercolor{subtitle}{fg=primary-gold}
  \setbeamercolor{author}{fg=primary-blue}
  \setbeamercolor{institute}{fg=neutral}
  \setbeamercolor{date}{fg=primary-gold}
  \setbeamercolor{itemize item}{fg=primary-blue}
  \setbeamercolor{itemize subitem}{fg=primary-gold}
  \setbeamercolor{alerted text}{fg=primary-gold}
  \setbeamercolor{block title}{fg=white, bg=primary-blue}
  \setbeamercolor{block body}{bg=light-bg}
  % Semantic block variants: block=definition/concept (blue), exampleblock=
  % worked example (green/positive), alertblock=key takeaway or caution (gold).
  % Keep to <=2 colored boxes per slide across ALL three variants (INV-7).
  \setbeamercolor{block title example}{fg=white, bg=positive}
  \setbeamercolor{block body example}{bg=light-bg}
  \setbeamercolor{block title alerted}{fg=white, bg=primary-gold}
  \setbeamercolor{block body alerted}{bg=light-bg}

  % Minimal footer: course code (left) + page X / N (right), no navigation clutter
  \setbeamertemplate{navigation symbols}{}
  \setbeamertemplate{footline}{%
    \color{neutral}\footnotesize\hspace{0.7em}\@coursecodetext\hfill
    \insertframenumber/\inserttotalframenumber\hspace{0.7em}\vspace{0.3em}}
}{}
\makeatother

% -----------------------------------------------------------------------------
% TikZ setup — libraries the snippet gallery uses
% -----------------------------------------------------------------------------
\usepackage{tikz}
\usetikzlibrary{arrows.meta, positioning, calc, decorations.pathreplacing,
                fit, shapes.geometric, backgrounds}

% Shared TikZ styles — snippets and hand-written diagrams can pick these up
% via `\begin{tikzpicture}[every node/.style={...}]` or by naming specific
% styles (e.g., `\node[dag-node] (X) at (0,0) {X};`).
\tikzset{
  % Node styles for causal DAGs and similar
  dag-node/.style = {
    draw=primary-blue, fill=light-bg, rounded corners=2pt,
    minimum width=1.2cm, minimum height=0.9cm, align=center, font=\small
  },
  decision-node/.style = {
    draw=primary-gold, fill=white, diamond, aspect=1.8,
    minimum width=1.8cm, minimum height=1.2cm, align=center, font=\small,
    inner sep=1pt
  },
  % Edge styles — observed vs counterfactual
  observed-edge/.style = {-{Latex[length=2.5mm]}, thick, draw=jet},
  counterfactual-edge/.style = {-{Latex[length=2.5mm]}, thick, draw=neutral, dashed},
  confound-edge/.style = {-{Latex[length=2.5mm]}, thick, draw=negative, dashed},
  % Data-point styles for plots
  observed-dot/.style = {fill=primary-blue, draw=primary-blue, circle, inner sep=1.2pt},
  counterfactual-dot/.style = {fill=white, draw=neutral, circle, inner sep=1.2pt}
}

% -----------------------------------------------------------------------------
% Convenience macros
% -----------------------------------------------------------------------------
\newcommand{\muted}[1]{\textcolor{neutral}{#1}}
\newcommand{\key}[1]{\textcolor{primary-gold}{\textbf{#1}}}
\newcommand{\good}[1]{\textcolor{positive}{#1}}
\newcommand{\bad}[1]{\textcolor{negative}{#1}}

% Standout frame macro: full-bleed dark background for section transitions.
% Beamer-only (uses \begin{frame}/\setbeamercolor/\usebeamerfont) -- do not
% call from an article-class file (e.g. Notes/<CODE>/*-notes.tex). \muted,
% \key, \good, \bad, and \coursecode above are plain xcolor/gdef and are
% safe to use from either class; verified when Notes/article-class support
% was added.
% Expand: \transitionslide{Your transition title}
\newcommand{\transitionslide}[1]{%
  \begingroup
  \setbeamercolor{background canvas}{bg=primary-blue}
  \begin{frame}[plain]
    \centering
    \vfill
    {\usebeamerfont{title}\color{white}\Large #1\par}
    \vfill
  \end{frame}
  \endgroup
}

% Section divider frame (Beamer-only), an alternative to \transitionslide with a
% darker two-line style: near-black (jet) background, a small white label line
% above a larger gold title, and a thin gold rule along the bottom edge.
% Expand: \sectiondivider{Section 2}{Pointers}
\newcommand{\sectiondivider}[2]{%
  \begingroup
  \setbeamercolor{background canvas}{bg=jet}
  \begin{frame}[plain]
    \vspace*{3.0cm}
    {\color{white}\ttfamily\large #1\par}
    \vspace{0.5cm}
    {\color{primary-gold}\LARGE #2\par}
    \begin{tikzpicture}[remember picture, overlay]
      \draw[primary-gold, line width=2.5pt]
        ($(current page.south west)+(0,0.1)$) -- ($(current page.south east)+(0,0.1)$);
    \end{tikzpicture}
  \end{frame}
  \endgroup
}

% -----------------------------------------------------------------------------
% Channel watermark — "Concepts That Click" (YouTube).
%
% Article-class files (CompetitiveExam Q/A sets, Notes, Assignments) get a
% light diagonal draft watermark on every page plus a centered footer naming
% the channel. Beamer decks get a subtle TikZ background watermark instead
% (the footline already carries the course code). Centralized here so every
% MCQ PDF is branded without per-file edits.
% -----------------------------------------------------------------------------
\makeatletter
\@ifclassloaded{article}{%
  \usepackage{draftwatermark}
  \SetWatermarkText{Concepts That Click}
  \SetWatermarkScale{1}
  \SetWatermarkFontSize{2cm}
  \SetWatermarkLightness{0.86}
  \SetWatermarkAngle{45}
  \usepackage{fancyhdr}
  \pagestyle{fancy}
  \fancyhf{}
  \fancyfoot[C]{\footnotesize\textcolor{neutral}{Concepts That Click~|~YouTube: \href{https://www.youtube.com/@concepts.thatclick}{@concepts.thatclick}}}
  \renewcommand{\headrulewidth}{0pt}
  \renewcommand{\footrulewidth}{0pt}
  \fancypagestyle{plain}{%
    \fancyhf{}%
    \fancyfoot[C]{\footnotesize\textcolor{neutral}{Concepts That Click~|~YouTube: \href{https://www.youtube.com/@concepts.thatclick}{@concepts.thatclick}}}%
    \renewcommand{\headrulewidth}{0pt}%
    \renewcommand{\footrulewidth}{0pt}%
  }
}{}
\@ifclassloaded{beamer}{%
  \setbeamertemplate{background}{%
    \begin{tikzpicture}[remember picture, overlay]
      \node[rotate=45, opacity=0.055, align=center]
        at (current page.center)
        {\Huge\textcolor{neutral}{Concepts That Click}};
    \end{tikzpicture}%
  }
}{}
\makeatother 
\coursecode{CS301}
\renewcommand{\thesection}{10.\arabic{section}}
\setlength{\emergencystretch}{3em}

\title{Coding Assignment 10: Hash Tables --- Chaining and Linear Probing}
\author{Auqib Hamid Lone\\[0.3em]
  \emph{Department of Computer Science \& Engineering}, \emph{GCET Ganderbal}}
\date{\today}

\begin{document}
\maketitle

\noindent\textbf{Due:} two weeks from issue. There is no automatic late penalty --- if you need more time, ask before the deadline and an extension will be granted (see the course policy in the syllabus).

\noindent This assignment autogrades Week~10's institute practical: \key{hash-table construction and collision handling}. The contract is given to you in \texttt{hash.h} --- \emph{do not modify it}. Copy \texttt{starter\_hash.c} to \texttt{hash.c} and fill in the functions. The setting is the running calculator. In Week~9 its $2{,}000$-identifier symbol table was sorted once and searched in $\Theta(\log n)$; this week the same names are found in one probe, and the price is that they are no longer in order. You build the table twice --- once with each collision policy from the deck --- and you must make both survive deletion.

\section{Why this problem}

Week~10's argument is that a lookup is not a search but an address computation: $\text{slot} = h(\text{key})$, with the \key{load factor} $\alpha = n/m$ as the one number that prices it. The deck shows the bill for $2{,}000$ identifiers --- about $1{,}000$ probes as a Week~6 list, $11$ as a Week~8 sorted array, $1.25$ as a hash table at $\alpha = 0.5$ --- and then spends its remaining time on the two things that keep the third number true. The first is the collision policy: \key{open hashing} (\key{separate chaining}), where colliding keys go in a per-slot list, versus \key{closed hashing} (\key{open addressing}), where they probe other slots in the same array. The second is the detail that quietly breaks student implementations: in open addressing, \emph{deletion cannot clear a slot}. An empty slot is the search's stopping condition, so clearing one hides every key that probed past it.

This assignment is both policies as running code, behind one contract, with the resize thresholds and the tombstone rule as explicit requirements. It is also the week's warning made executable: a hash table that never rehashes is still ``correct'' on small tests and still $O(n)$ on real ones, so the hidden suite measures both.

\section{Your task}

Implement the functions declared in \texttt{hash.h}. There are three groups.

\textbf{The two hash functions.} \texttt{hash\_string} is the deck's polynomial (Horner) scheme: base $B = 31$, treated as unsigned, reduced modulo $m$ \emph{at every step} --- the reduction inside the loop is what stops $31^{\text{len}}$ from overflowing. \texttt{hash\_int} is the division method, $h(k) = k \bmod m$, with the key read as unsigned so that negative keys stay in range. Both are stated exactly in the header; do not ``improve'' them, because the tests check exact values.

\textbf{Open hashing --- separate chaining (\texttt{ChainTable}).}
\begin{verbatim}
ChainTable *chain_create(unsigned m);
int   chain_insert(ChainTable *t, const char *key, int value);
int   chain_search(const ChainTable *t, const char *key, int *out);
int   chain_delete(ChainTable *t, const char *key);
int   chain_count(const ChainTable *t);
unsigned chain_slots(const ChainTable *t);
double   chain_load_factor(const ChainTable *t);
void  chain_destroy(ChainTable *t);
\end{verbatim}
Slot $j$ holds the head of a list of every key with $h(\text{key}) = j$. Insert pushes a new node, search walks that one list, and delete unlinks a node --- no tombstones, because a chain is a real list. $\alpha$ may exceed $1$, but you must not let it: if inserting a new key would make $n > m$, rehash into a larger slot array first, so that after any insert $n \le m$.

\textbf{Closed hashing --- open addressing with linear probing (\texttt{OpenTable}).}
\begin{verbatim}
OpenTable *oa_create(unsigned m);
int   oa_insert(OpenTable *t, const char *key, int value);
int   oa_search(const OpenTable *t, const char *key, int *out);
int   oa_delete(OpenTable *t, const char *key);
int   oa_count(const OpenTable *t);
int   oa_tombstones(const OpenTable *t);
unsigned oa_slots(const OpenTable *t);
double   oa_load_factor(const OpenTable *t);
void  oa_destroy(OpenTable *t);
\end{verbatim}
Every key lives inside the table. The probe sequence is $h, h+1, h+2, \dots \pmod m$ with $h = \texttt{hash\_string}(\text{key}, m)$. Each slot has \emph{three} states, not two: \texttt{EMPTY}, \texttt{OCCUPIED}, \texttt{DELETED}. Search stops only at \texttt{EMPTY}; it walks past \texttt{DELETED}. Insert may reuse the first \texttt{DELETED} slot it passes. Delete must mark the slot \texttt{DELETED} (a \key{tombstone}) and increment the tombstone count --- never clear it.

Two thresholds keep linear probing from degrading. After any insert that returns $1$ (a genuinely new key), both must hold: occupancy $n + \text{tombstones} \le m/2$, and $\text{tombstones} \le m/4$. Grow the slot array when the first would be broken; rebuild the table (same $m$ is enough) when the second would be.

\textbf{Policy shared by both tables.} A key is unique: inserting a key that is already present \emph{replaces} its value, leaves the count unchanged, and returns $0$; a new key returns $1$. The table stores its own \emph{copy} of every key --- the caller's buffer may be a stack array that is overwritten immediately --- and releases every copy in delete, rehash and destroy. Return values, NULL-safety and failure cases are stated function by function in \texttt{hash.h}.

\section{Constraints}

\begin{itemize}
  \item Compile with \texttt{gcc -std=c11 -Wall -Wextra -Werror}: no errors and no warnings. No global or static mutable state.
  \item \texttt{m == 0} is rejected: \texttt{create} returns \texttt{NULL}. Every other function is NULL-safe on a NULL table, as documented in the header.
  \item Chaining: after any insert that returns $1$, \texttt{chain\_count(t) <= (int)chain\_slots(t)}; the expected cost of insert, search and delete is $O(1+\alpha)$.
  \item Open addressing: after any insert that returns $1$, $2n \le m$ and $\text{tombstones} \le m/4$; deletion must leave a tombstone; the expected probe count is $O(1/(1-\alpha))$.
  \item Neither table may degrade to $O(n)$ per operation: the hidden suite runs $20{,}000$ keys through each and imposes a wall-clock timeout.
  \item Keys are arbitrary C strings, including \texttt{""}; values are \texttt{int} and may be negative. Rehashing must move key copies, not duplicate or leak them.
  \item The whole suite runs under \texttt{-fsanitize=address,undefined}: a leak, a double free or undefined behaviour is a failure, not a deduction.
\end{itemize}

\section{Worked examples (these are the visible tests)}

\textbf{Example 1 --- the hash functions.} With $m = 13$: \texttt{hash\_int(40,13) = 1}, \texttt{hash\_int(53,13) = 1}, \texttt{hash\_int(66,13) = 1}, \texttt{hash\_int(100,13) = 9}, and \texttt{hash\_int(-1,13) = 8} (the key is read as unsigned). With $m = 101$: \texttt{hash\_string("dog",101) = 58} and \texttt{hash\_string("",101) = 0}. These are exact values, not properties --- check them before building anything on top.

\textbf{Example 2 --- separate chaining.} \texttt{m = 5}, so slot $=$ \texttt{hash\_string(key, 5)}: alpha $\to 3$, beta $\to 2$, gamma $\to 0$, delta $\to 2$.

\begin{verbatim}
insert alpha=10, beta=20, gamma=30, delta=40   (each returns 1)
  slot 0: gamma      slot 1: --      slot 2: delta -> beta
  slot 3: alpha      slot 4: --
count = 4, load factor = 0.8
search delta  -> 1 (value 40)      search beta -> 1 (value 20)
insert beta=25 -> 0 (replaced), count still 4, search beta -> 25
delete alpha  -> 1;  search alpha -> 0;  search delta -> 1
insert alpha=99 -> 1;  search alpha -> 99
\end{verbatim}
Note the last line: delete then reinsert is an insert, not a replacement, so it returns $1$. Beta and delta share slot $2$; deleting one must not disturb the other.

\textbf{Example 3 --- linear probing and the tombstone.} \texttt{m = 11}, so alpha $\to 2$, beta $\to 2$, gamma $\to 10$, delta $\to 0$. Inserting alpha, beta, gamma, delta in that order places them at slots $2, 3, 10, 0$: beta's home slot is taken, so it takes the next one. Four live entries in eleven slots keeps occupancy at $8 \le 11$, so no rehash occurs and the trace stays visible.

\begin{verbatim}
delete alpha        -> 1;  slot 2 becomes DELETED; tombstones = 1
search beta         -> 1;  probes 2 (DELETED, keep going) then 3 -> found
search alpha        -> 0
insert alpha=99     -> 1;  reuses slot 2;  tombstones back to 0
\end{verbatim}
Had \texttt{delete} cleared slot $2$ instead, \texttt{search(beta)} would stop at the empty slot and report beta absent --- while beta sits one slot later. That single mistake is what the tombstone checks look for.

\textbf{Example 4 --- \texttt{m = 1}.} The smallest legal table. Twelve inserts into \texttt{chain\_create(1)} and \texttt{oa\_create(1)} must all succeed, all twelve keys must be found, and both tables must have grown past one slot. A one-slot table is where an off-by-one in the resize condition shows up.

\textbf{Example 5 --- teardown.} \texttt{destroy(NULL)} is safe; \texttt{chain\_create(0)} and \texttt{oa\_create(0)} return \texttt{NULL}. Every table you create must be destroyed; the sanitizer reports the keys and nodes you forget.

\section{How to work}

\begin{enumerate}
  \item Copy \texttt{starter\_hash.c} to \texttt{hash.c}; keep \texttt{hash.h} unchanged.
  \item Build and run the visible tests:
\begin{verbatim}
gcc -std=c11 -Wall -Wextra -Werror -o test_visible test_visible.c hash.c
./test_visible
\end{verbatim}
  \item Get Examples 1--2 green before writing any probing code. A wrong hash function makes every later test look like a collision bug.
  \item Trace Example 3 on paper twice: once with a tombstone, once with a cleared slot. Then write the probe loop so that \texttt{EMPTY} is the only stopping condition.
  \item Run the same build under the sanitizers before submitting:
\begin{verbatim}
gcc -std=c11 -Wall -Wextra -Werror -g \
    -fsanitize=address,undefined -o test_visible test_visible.c hash.c
./test_visible
\end{verbatim}
  \item The hidden suite adds: randomized insert/search/delete against a reference model (duplicates included), forced single-bucket collisions, deletion of the earliest key of a probe run, delete-then-reinsert, tombstone-bound and occupancy-bound checks, $20{,}000$-key scale runs for both tables, and leak checking.
\end{enumerate}

\section{Submission and grading}

Submit \texttt{hash.c} through the course submission form. Grading is by the autograder --- correctness on the hidden suite, compiled with \texttt{-Werror} and run under \texttt{-fsanitize=address,undefined} --- plus a short code read for the four structural requirements: a real per-slot chain, a real tombstone on delete, an occupancy/tombstone rehash that keeps both bounds, and key copies that are owned and released. A table that is correct on the small cases but never rehashes is partially credited; a warning under \texttt{-Werror}, a leak, or a sanitizer report is a failure.

\bibliography{../../Bibliography_base}
\bibliographystyle{plain}
\end{document}
